\SourceLeaf{222}%
von $h,a,r$ für $\rho=\rho'$ mit $h',a',r'$ bezeichnet sind. Als Constante hat man den Werth von $\frac{h(a-x)}{r}$ für $\rho=0$ anzunehmen, welcher wenn man die Dichtigkeit in $P$ mit $k^0$ bezeichnet, $=-k^0$ wird für ein positives $x$, und $=+k^0$ für ein negatives, indem für $\rho=0$ offenbar $a=0$, $\psi=0$, $h=k^0$, $x=\pm r$ wird. Für den Fall $x=0$ hingegen hat man als Constante den Grenzwerth von $\frac{ha}{r}$ bei unendlich abnehmendem $\rho$ anzunehmen, welcher $=0$ ist, weil $a$ ein Unendlichkleines von einer höhern Ordnung wird als $r$.

Der Werth des Integrals $\int\frac{\mathrm dh}{\mathrm d\rho}.\frac{a-x}{r}.\mathrm d\rho$ bleibt bis auf einen unendlich kleinen Unterschied derselbe, man möge $x=0$, oder unendlich klein $=\pm\varepsilon$ setzen. Zerlegt man nemlich jenes Integral in
\[
\int_0^\delta\frac{\mathrm dh}{\mathrm d\rho}.\frac{a-x}{r}.\mathrm d\rho+\int_\delta^{\rho'}\frac{\mathrm dh}{\mathrm d\rho}.\frac{a-x}{r}.\mathrm d\rho
\]
so ist klar, dass das Behauptete für den ersten Theil gilt, wenn $\delta$ unendlich klein, und für den zweiten, wenn $\frac{\delta}{\varepsilon}$ unendlich gross ist, also für das Ganze, wenn $\delta$ ein Unendlichkleines von einer niedrigern Ordnung als $\varepsilon$.

Ein ähnlicher Schluss gilt auch in Beziehung auf das Integral $\int\frac{\mathrm da}{\mathrm d\rho}.\frac{h\rho\rho}{r^3}.\mathrm d\rho$, wenn die Punkte der Fläche, welche dem bestimmten Werthe von $\theta$ entsprechen, eine Curve bilden, die in $P$ eine messbare Krümmung hat, so dass $\frac{a}{\rho\rho}$ in dem hier betrachteten Raume einen endlichen nach der Stetigkeit sich ändernden Werth erhält. Bezeichnet man nemlich diesen Werth mit $A$, so wird
\[
\frac{\mathrm da}{\mathrm d\rho}=2A\rho+\frac{\mathrm dA}{\mathrm d\rho}.\rho\rho
\]
mithin zerlegt sich jenes Integral in folgende zwei
\[
\int\frac{2\rho^3Ah\mathrm d\rho}{r^3}+\int\frac{\mathrm dA}{\mathrm d\rho}.\frac{\rho^4}{r^3}h\mathrm d\rho
\]
bei welchen beiden die Gültigkeit obiger Schlussweise von selbst klar ist.

Endlich sind auch offenbar die Werthe von $\frac{h'(a'-x)}{r'}$ für alle drei Werthe von $x$ bis auf unendlich kleine Unterschiede gleich.

Hieraus folgt also, dass $Q'+k^0$, $Q^0$, $Q''-k^0$ bis auf unendlich kleine Unterschiede gleich sind, und dasselbe wird demnach auch von
\[
\int(Q'+k^0)\mathrm d\theta,\quad\int Q^0\mathrm d\theta,\quad\int(Q''-k^0)\mathrm d\theta
\]
gelten, oder von den Grössen
\[
X'+2\pi k^0,\quad X^0,\quad X''-2\pi k^0
\]

\SourceLeaf{223}%
Man kann diesen wichtigen Satz auch so ausdrücken: der Grenzwerth von $X$, bei unendlich abnehmendem positiven $x$ ist $X^0-2\pi k^0$, bei unendlich abnehmendem negativen $x$ hingegen $X^0+2\pi k^0$, oder $X$ ändert sich zweimal sprungsweise um $-2\pi k^0$, indem $x$ aus einem negativen Werthe in einen positiven übergeht, das erstemal, indem $x$ den Werth $0$ erreicht, und das zweitemal, indem es ihn überschreitet.

\Article{16}
In der Beweisführung des vorhergehenden Artikels ist zwar vorausgesetzt, dass die Schnitte der Fläche mit den durch die erste Coordinatenaxe gelegten Ebenen in $P$ eine messbare Krümmung haben: allein unser Resultat bleibt auch noch gültig, wenn die Krümmung in $P$ unendlich gross ist, einen einzigen Fall ausgenommen. Dass $\frac a\rho$ für ein unendlich kleines $\rho$ selbst unendlich klein werden müsse, bringt schon die Voraussetzung des Vorhandenseins einer bestimmten Berührungsebene an der Fläche in $P$ mit sich; allein von einerlei Ordnung sind beide Grössen nur dann, wenn ein endlicher Krümmungshalbmesser Statt findet; bei einem unendlich kleinen Krümmungshalbmesser hingegen wird $\frac a\rho$ von einer niedrigern Ordnung sein, als $\rho$. Wir werden nun zeigen, dass unsre Resultate auch im letztern Falle ihre Gültigkeit behalten, wenn nur die Ordnungen beider Grössen \textit{vergleichbar} sind.

Nehmen wir also an, $\frac a\rho$ sei von derselben Ordnung wie $\rho^\mu$, wo $\mu$ einen endlichen positiven Exponenten bedeutet, also $\frac{a}{\rho^{1+\mu}}$ eine endliche in dem in Rede stehenden Raume nach der Stetigkeit sich ändernde Grösse, die wir mit $B$ bezeichnen wollen. Es zerfällt also das Integral $\int\frac{\mathrm da}{\mathrm d\rho}.\frac{h\rho\rho}{r^3}\mathrm d\rho$ in die beiden folgenden
\[
\int\frac{(1+\mu)\rho^{2+\mu}hB\mathrm d\rho}{r^3}+\int\frac{\rho^{3+\mu}}{r^3}.\frac{\mathrm dB}{\mathrm d\rho}.h\mathrm d\rho
\]
Auf das zweite Integral lassen sich die Schlüsse des vorhergehenden Artikels unmittelbar anwenden, auf das erste hingegen nach einer leichten Umformung. Setzt man nemlich $\frac1\mu=m$, $\rho^\mu=\sigma$ oder $\rho=\sigma^m$, so wird jenes Integral
\[
=(m+1)\int\frac{Bh\sigma^{3m}\mathrm d\sigma}{(\sigma^{2m}+(a-x)^2)^{\frac32}}
\]
Auch dieses Integral hat nun offenbar so lange nur einen unendlich kleinen Werth, als die Integration nur von $0$ bis zu einem unendlich kleinen Werthe von $\sigma$ ausgedehnt wird; für jeden endlichen Werth von $\sigma$ hingegen erhält der Coëfficient
\SourceLeaf{224}%
von $\mathrm d\sigma$ bis auf einen unendlich kleinen Unterschied einerlei Werth, man möge $x=0$ oder unendlich klein annehmen. Dies gilt also auch von dem ganzen Integral, wenn es von $\sigma=0$ bis $\sigma=\BareRoot[m]{\rho'}$ ausgedehnt wird.

Nur in einem einzigen Falle verlieren unsre Schlüsse ihre Gültigkeit, wenn nemlich $\frac a\rho$ mit keiner Potenz von $\rho$ mehr zu einerlei Ordnung gehört, wie z. B. wenn $\frac a\rho$ von derselben Ordnung wäre, wie $\frac{1}{\log\frac1\rho}$. In diesem Falle würde $Q$ bei unendlicher Annäherung des Punktes $O$ zur Fläche über alle Grenzen wachsen, und dasselbe würde auch für $X$ gelten, wenn ein solches Verhalten nicht bloss für einen oder einige Werthe von $\theta$, sondern für alle Statt fände. Es ist jedoch unnöthig, dies hier weiter zu entwickeln, da wir diesen singulären Fall von unsrer Untersuchung ohne Nachtheil ganz ausschliessen können.

\Article{17}
Wir wollen nun unter denselben Voraussetzungen und Bezeichnungen, wie im 15. Artikel, die Grösse $Y$ betrachten, wovon $\frac{hb\mathrm db.\mathrm dc}{r^3}$ ein unbestimmtes Element ist. Da $r=\BareRoot{(bb+cc+(a-x)^2)}$, und folglich
\[
\frac{\mathrm d\frac hr}{\mathrm db}=-\frac{hb}{r^3}+\frac1r.\frac{\mathrm dh}{\mathrm db}-\frac{h(a-x)}{r^3}.\frac{\mathrm da}{\mathrm db}
\]
insofern $c$ als constant betrachtet wird, so gibt die erste Integration in diesem Sinne,
\[
\int\frac{hb\mathrm db}{r^3}=\frac{h^*}{r^*}-\frac{h^{**}}{r^{**}}+\int\frac1r.\frac{\mathrm dh}{\mathrm db}.\mathrm db-\int\frac{h(a-x)}{r^3}.\frac{\mathrm da}{\mathrm db}.\mathrm db
\]
wo die Integrationen sich vom kleinsten zum grössten Werthe von $b$, für jeden bestimmten Werth von $c$ erstrecken, und mit $h^*,r^*,h^{**},r^{**}$ die jenen Grenzwerthen entsprechenden Werthe von $h$ und $r$ bezeichnet sind, Schreiben wir zur Abkürzung
\[
\frac{h^*}{r^*}-\frac{h^{**}}{r^{**}}=T,\quad\frac\rho r.\frac{\mathrm dh}{\mathrm db}-\frac{h(a-x)\rho}{r^3}.\frac{\mathrm da}{\mathrm db}=U
\]
so wird
\[
Y=\int T\mathrm dc+\iint\frac U\rho.\mathrm db.\mathrm dc
\]
\SourceLeaf{225}%
wo die Integration in Beziehung auf $c$ vom kleinsten Werthe, welchen diese Coordinate in der Fläche hat, bis zum grössten ausgedehnt werden muss. In dem doppelten Integrale stellt $\mathrm db.\mathrm dc$ die Projection eines unbestimmten Elements der Fläche auf die Ebene der $b,c$ vor, und es kann mithin auch $\rho\mathrm d\rho.\mathrm d\theta$ dafür geschrieben werden; sonach wird
\[
Y=\int T\mathrm dc+\iint U\mathrm d\rho.\mathrm d\theta
\]
wo in dem Doppelintegral von $\rho=0$ bis $\rho=\rho'$ und von $\theta=0$ bis $\theta=2\pi$ integrirt werden muss. Durch ähnliche Schlüsse, wie im 15. Artikel, erkennt man nun leicht, dass dieser Ausdruck bis auf unendlich kleine Unterschiede gleiche Werthe erhält, man möge $x=0$ oder unendlich klein annehmen, oder mit andern Worten, der Werth von $Y$ hat bei positiven und bei negativen unendlich abnehmenden Werthen von $x$ eine und dieselbe Grenze, und diese Grenze ist nichts anderes, als der Werth obiger Formel, wenn man darin $x=0$ setzt. Wir wollen nach der Analogie diesen Werth mit $Y^0$ bezeichnen, wobei jedoch bemerkt werden muss, dass man nicht sagen darf, es sei dies \textit{der} Werth von $\int\frac{kb\mathrm ds}{r^3}$ für $x=0$ (insofern dieser Ausdruck für $x=0$ eine wahre Integration nicht zulässt), sondern nur, es sei \textit{ein} Werth jenes Integrals, nemlich derjenige, welcher hervorgeht, wenn man in der oben befolgten Ordnung integrirt.

Übrigens bedarf dieses Resultat (auf ähnliche Weise wie oben Art. 16) einer Einschränkung in dem singulären Falle, wo in dem Punkte $P$ unendlich kleine Krümmungshalbmesser Statt finden, imgleichen, wenn in diesem Punkte $\frac{\mathrm dh}{\mathrm db}$ unendlich gross wird: für unsern Zweck ist es jedoch unnöthig, solche Ausnahmsfälle, die nur in einzelnen Punkten oder Linien vorkommen können (also nicht in Theilen der Fläche, sondern nur an der Grenze von Theilen), besonders zu betrachten.

Endlich ist von selbst klar, dass es sich mit der Grösse $Z$ oder dem Integrale $\int\frac{kc\mathrm ds}{r^3}$ ganz eben so verhält, wie mit $Y$, nemlich dass dieses Integral, wenn der Punkt $O$ sich in der ersten Coordinatenaxe dem Punkte $P$ unendlich nähert, einerlei Grenzwerth $Z^0$ hat, die Annäherung mag auf der positiven oder auf der negativen Seite Statt finden, und dass dieser Grenzwerth zugleich der Werth von $\iint\frac{hc\mathrm dc.\mathrm db}{r^3}$ für $x=0$ ist, insofern man zuerst nach $c$ integrirt.

\SourceLeaf{226}%
\Article{18}
Erwägen wir nun, dass die Grössen $\frac{\mathrm dV}{\mathrm dx}$, $\frac{\mathrm dV}{\mathrm dy}$, $\frac{\mathrm dV}{\mathrm dz}$ in allen Punkten des Raums, die nicht in der Fläche selbst liegen, unbedingt einerlei sind mit $X,Y,Z$, und dass $V$ sich überall nach der Stetigkeit ändert, so lässt sich aus den in dem vorhergehenden Artikel gefundenen Resultaten leicht folgern, dass in unendlich kleiner Entfernung von $P$, oder für unendlich kleine Werthe von $x,y,z$ der Werth von $V$ bis auf unendlich kleine Grössen höherer Ordnung genau, ausgedrückt wird durch
\[
V^0+x(X^0-2\pi k^0)+yY^0+zZ^0
\]
wenn $x$ positiv ist, oder durch
\[
V^0+x(X^0+2\pi k^0)+yY^0+zZ^0
\]
wenn $x$ negativ ist, wo mit $V^0$ der Werth von $V$ in dem Punkte $P$ selbst, oder für $x=0$, $y=0$, $z=0$ bezeichnet ist. Betrachten wir also die Werthe von $V$ in einer durch $P$ gelegten geraden Linie, die mit den drei Coordinataxen die Winkel $A,B,C$ macht, bezeichnen mit $t$ ein unbestimmtes Stück dieser Linie und mit $t^0$ den Werth von $t$ in dem Punkte $P$, so wird, wenn $t-t^0$ unendlich klein ist, bis auf ein Unendlichkleines höherer Ordnung genau
\[
V=V^0+(t-t^0)(X^0\cos A+Y^0\cos B+Z^0\cos C\mp2\pi k^0\cos A)
\]
das obere Zeichen für positive, das untere für negative Werthe von $(t-t^0)\cos A$ geltend, oder es hat $\frac{\mathrm dV}{\mathrm dt}$ in dem Punkte $P$ für ein spitzes $A$ zwei verschiedene Werthe, nemlich
\[
\begin{gathered}
X^0\cos A+Y^0\cos B+Z^0\cos C-2\pi k^0\cos A\text{ und}\\
X^0\cos A+Y^0\cos B+Z^0\cos C+2\pi k^0\cos A
\end{gathered}
\]
je nachdem $\mathrm dt$ als positiv oder als negativ betrachtet wird. Für den Fall, wo $A$ ein rechter Winkel ist, also die gerade Linie die Fläche nur berührt, fallen beide Ausdrücke zusammen, und es wird
\[
\frac{\mathrm dV}{\mathrm dt}=Y^0\cos B+Z^0\cos C
\]
\begin{center}*\qquad\qquad*\\{}*\end{center}

\SourceLeaf{227}%
Die bisher vorgetragenen Sätze sind zwar ihrem wesentlichen Inhalte nach nicht neu, durften aber des Zusammenhanges wegen. als nothwendige Vorbereitungen zu den nachfolgenden Untersuchungen nicht übergangen werden, in. welchen eine Reihe neuer Lehrsätze entwickelt werden wird.

\Article{19}
Es sei $V$ das Potential eines Systems von Massen $M',M'',M'''\ldots$, die sich in den Punkten $P',P'',P'''\ldots$ befinden; $v$ das Potential eines zweiten Systems von Massen $m',m'',m'''\ldots$, die in den Punkten $p',p'',p'''\ldots$ angenommen werden; ferner seien $V',V'',V'''\ldots$ die Werthe von $V$ in den letztern Punkten, und $v',v'',v'''\ldots$ die Werthe von $v$ in den Punkten $P',P'',P'''\ldots$ Man hat dann die Gleichung
\[
M'v'+M''v''+M'''v'''+\text{u. s. f.}=m'V'+m''V''+m'''V'''+\text{u. s. f.}
\]
die auch durch $\sum Mv=\sum mV$ ausgedrückt wird, wenn unbestimmt $M$ jede Masse des ersten, $m$ jede Masse des zweiten Systems vorstellt. In der That ist sowohl $\sum Mv$ als $\sum mV$ nichts anderes, als das Aggregat aller Combinationen $\frac{Mm}{\rho}$, wenn $\rho$ die gegenseitige Entfernung der Punkte bezeichnet, in welchen sich die betreffenden Massen $M,m$ befinden.

Befinden sich die Massen des einen Systems, oder beider, nicht in discreten Punkten, sondern auf Linien, Flächen oder körperliche Räume nach der Stetigkeit vertheilt, so behält obige Gleichung ihre Gültigkeit, wenn man anstatt der Summe das entsprechende Integral substituirt.

Ist also z. B. das zweite Massensystem in einer Fläche so vertheilt, dass auf das Flächenelement $\mathrm ds$ die Masse $k\mathrm ds$ kommt, so wird $\sum Mv=\int kV\mathrm ds$, oder wenn ähnliches auch von dem ersten System gilt, so dass das Flächenelement $\mathrm dS$ die Masse $K\mathrm dS$ enthält, wird $\int Kv\mathrm dS=\int kV\mathrm ds$. Es ist von Wichtigkeit, in Beziehung auf letztern Fall zu bemerken, dass diese Gleichung noch gültig bleibt, wenn beide Flächen coincidiren; der Kürze wegen wollen wir aber die Art, wie diese Erweiterung des Satzes strenge gerechtfertigt werden kann, hier jetzt nur nach ihren Hauptmomenten andeuten. Es ist nemlich nicht schwer nachzuweisen, dass diese beiden Integrale, insofern sie sich auf Eine und dieselbe Fläche beziehen, die Grenzwerthe von denen sind, die sich auf zwei getrennte Flächen
\SourceLeaf{228}%
beziehen, indem man die Entfernung derselben von einander unendlich abnehmen lässt, zu welchem Zweck man nur diese beiden Flächen gleich und parallel anzunehmen braucht. Unmittelbar einleuchtend ist zwar diese Beweisart nur in sofern, als die vorgegebene Fläche so beschaffen ist, dass die Normalen in allen ihren Punkten mit Einer geraden Linie spitze Winkel machen. Eine Fläche, wo diese Bedingung fehlt (wie allemal, wenn von einer geschlossenen Fläche die Rede ist), wird zuvor in zwei oder mehrere Theile zu zerlegen sein, die einzeln jener Bedingung Genüge leisten, wodurch es leicht wird, diesen Fall auf den vorigen zurückzuführen.

\Article{20}
Wenden wir das Theorem des vorhergehenden Artikels auf den Fall an, wo das zweite Massensystem mit gleichförmiger Dichtigkeit $k=1$ auf eine Kugelfläche vertheilt ist, deren Halbmesser $=R$, so ist das daraus entspringende Potential $v$ im Innern der Kugel constant $=4\pi R$; in jedem Punkte ausserhalb der Kugel, dessen Entfernung vom Mittelpunkte $=r$, wird $v=\frac{4\pi RR}{r}$, oder eben so gross, wie im Mittelpunkte das Potential von einer in jenem Punkte angenommenen Masse $4\pi RR$; auf der Oberfläche der Kugel fallen beide Werthe von $v$ zusammen. Befindet sich also das erste Massensystem ganz im Innern der Kugel, so wird $\sum Mv$ äqual dem Producte der Gesammtmasse dieses Systems in $4\pi R$; ist aber jenes Massensystem ganz ausserhalb der Kugel, so wird $\sum Mv$ äqual dem Producte des Potentials dieser Masse im Mittelpunkte der Kugel in $4\pi RR$; ist endlich das erste Massensystem auf der Oberfläche der Kugel nach der Stetigkeit vertheilt, so sind für $\int Kv\mathrm dS$ beide Ausdrücke gleichgültig. Es folgt hieraus der

\textsc{Lehrsatz.} Bedeutet $V$ das Potential einer wie immer vertheilten Masse in dem Elemente einer mit dem Halbmesser $R$ beschriebenen Kugelfläche $\mathrm ds$, so wird, durch die ganze Kugelfläche integrirt,
\[
\int V\mathrm ds=4\pi(RM^0+RRV^0)
\]
wenn man mit $M^0$ die ganze im Innern der Kugel befindliche Masse, mit $V^0$ das Potential der ausserhalb befindlichen Masse im Mittelpunkt der Kugel bezeichnet, und dabei die Massen, die etwa auf der Oberfläche der Kugel stetig vertheilt sein mögen, nach Belieben den äussern oder innern Massen zuordnet.

\SourceLeaf{229}%
\Article{21}
\textsc{Lehrsatz.} Das Potential $V$ von Massen, die sämmtlich ausserhalb eines zusammenhängenden Raumes liegen, kann nicht in einem Theile dieses Raumes einen constanten Werth und zugleich in einem andern Theile desselben einen verschiedenen Werth haben.

\textit{Beweis.} Nehmen wir an, es sei in jedem Punkte des Raums $A$ das Potential constant $=a$, und in jedem Punkte eines andern an $A$ grenzenden keine Masse enthaltenden Raumes $B$ (algebraisch) grösser als $a$. Man construire eine Kugel, wovon ein Theil in $B$, der übrige Theil aber nebst dem Mittelpunkte in $A$ enthalten ist, welche Construction allemal möglich sein wird. Ist nun $R$ der Halbmesser dieser Kugel, und $\mathrm ds$ ein unbestimmtes Element ihrer Oberfläche, so ist nach dem Lehrsatze des vorigen Artikels $\int V\mathrm ds=4\pi RRa$, und $\int(V-a)\mathrm ds=0$, was unmöglich ist, da für den Theil der Oberfläche, welcher in $A$ liegt, $V-a=0$, und für den übrigen Theil der Voraussetzung zu Folge nicht $=0$, sondern positiv ist.

Auf ganz ähnliche Weise erhellt die Unmöglichkeit, dass in allen Punkten eines an $A$ grenzenden Raumes $V$ kleiner sei als $a$.

Offenbar müsste aber wenigstens einer dieser beiden Fälle Statt finden, wenn unser Theorem falsch wäre.

Dieser Lehrsatz enthält folgende zwei Sätze:

I. Wenn der die Massen enthaltende Raum schalenförmig einen massenleeren Raum umschliesst, und das Potential in einem Theile dieses Raumes einen constanten Werth hat, so gilt dieser für alle Punkte des ganzen eingeschlossenen Raumes.

II. Wenn das Potential der in einen endlichen Raum eingeschlossenen Massen in irgend einem Theile des äussern Raumes einen constanten Werth hat, so gilt dieser für den ganzen unendlichen äussern Raum.

Zugleich erhellt leicht, dass in diesem zweiten Fall der constante Werth des Potentials kein anderer als $0$ sein kann. Denn wenn man mit $M$ das Aggregat aller Massen, falls sie sämmtlich einerlei Zeichen haben, oder im entgegengesetzten Fall das Aggregat der positiven oder der negativen Massen allein, je nachdem jene oder diese überwiegen, bezeichnet, so ist das Potential in einem Punkte, dessen Entfernung von dem nächsten Massenelemente $=r$ ist, jedenfalls, abso%
\SourceLeaf{230}%
lut genommen, kleiner als $\frac Mr$, welcher Bruch offenbar im äussern Raume kleiner als jede angebliche Grösse werden kann.

\Article{22}
\textsc{Lehrsatz.} Ist $\mathrm ds$ das Element einer einen zusammenhängenden endlichen Raum begrenzenden Fläche, $P$ die Kraft, welche irgendwie vertheilte Massen in $\mathrm ds$ in der auf die Fläche normalen Richtung ausüben, wobei eine nach innen oder nach aussen gerichtete Kraft als positiv betrachtet wird, je nachdem anziehende oder abstossende Massen als positiv gelten: so wird das Integral $\int P\mathrm ds$ über die ganze Fläche ausgedehnt $=4\pi M+2\pi M'$, wenn $M$ das Aggregat der im Innern des Raumes befindlichen, $M'$ das der auf der Oberfläche nach der Stetigkeit vertheilten Massen bedeuten.

\textit{Beweis.} Bezeichnet man mit $U\mathrm d\mu$ denjenigen Theil von $P$, welcher von dem Massenelemente $\mathrm d\mu$ herrührt, mit $r$ die Entfernung des Elements $\mathrm d\mu$ von $\mathrm ds$, und mit $u$ den Winkel, welchen in $\mathrm ds$ die nach Innen gerichtete Normale mit $r$ macht, so ist $U=\frac{\cos u}{rr}$. Es ist aber in Beziehung auf jedes bestimmte $\mathrm d\mu$, vermöge eines in der \textit{Theoria Attractionis corporum sphaeroidicorum ellipticorum} Art. 6 bewiesenen Lehrsatzes $\int\frac{\cos u}{rr}.\mathrm ds=0$, $2\pi$ oder $4\pi$, je nachdem $\mathrm d\mu$ ausserhalb des durch die Fläche begrenzten Raumes, in der Fläche selbst, oder innerhalb jenes Raumes liegt. Da nun $\int P\mathrm ds$ dem Gesammtbetrage aller $\mathrm d\mu.\int U\mathrm ds$ gleichkommt, so ergibt sich hieraus unser Theorem von selbst.

In Beziehung auf den hier benutzten Hülfssatz muss noch bemerkt werden, dass derselbe, in der Gestalt wie er a. a. O. ausgesprochen ist, für einen speciellen Fall einer Modification bedarf. Es bedeutet nemlich $r$ die Entfernung eines \textit{gegebenen Punktes} von dem Elemente $\mathrm ds$, und für den Fall, wo dieser Punkt in der Fläche selbst liegt, ist die Formel $\int\frac{\cos u}{rr}.\mathrm ds=2\pi$ nur insofern richtig, als die Stetigkeit der Krümmung der Fläche in dem Punkte nicht verletzt wird. Eine solche Verletzung findet aber Statt, wenn der Punkt in einer Kante oder Ecke liegt, und dann muss anstatt $2\pi$ der Inhalt derjenigen Figur gesetzt werden, welche durch die sämmtlichen von da ausgehenden die Fläche tangirenden geraden Linien aus einer um den Punkt als Mittelpunkt mit dem Halbmesser $1$ beschriebenen Kugelfläche ausgeschieden wird. Da jedoch solche Ausnahmsfälle nur Linien oder Punkte, also nicht \textit{Theile} der Fläche, sondern nur Scheidungs%
\SourceLeaf{231}%
grenzen zwischen Theilen betreffen, so hat dies offenbar auf die von dem Hülfssatze hier gemachte Anwendung gar keinen Einfluss.

\Article{23}
Wir legen durch jeden Punkt der Fläche eine Normale, und bezeichnen mit $p$ die Entfernung eines unbestimmten Punktes derselben von dem in die Fläche selbst gesetzten Anfangspunkte, auf der innern Seite der Fläche als positiv betrachtet. Das Potential der Massen $V$ kann als Function von $p$ und zweien andern veränderlichen Grössen betrachtet werden, die auf irgendwelche Art die einzelnen Punkte der Fläche von einander unterscheiden, und eben so verhält es sich mit dem partiellen Differentialquotienten $\frac{\mathrm dV}{\mathrm dp}$, dessen Werth hier aber nur für die in die Fläche selbst fallenden Punkte, oder für $p=0$ in Betracht gezogen werden soll. Da dieser mit $P$ völlig gleichbedeutend ist, wenn Massen sich nur in dem innern Raume, oder in dem äussern, oder in beiden befinden, keine Masse aber auf die Fläche selbst vertheilt ist, so hat man in diesem Falle
\[
\int\frac{\mathrm dV}{\mathrm dp}.\mathrm ds=4\pi M
\]
In dem Falle hingegen, wo die ganze Masse bloss auf der Fläche selbst vertheilt ist, so dass das Element $\mathrm ds$ die Masse $k\mathrm ds$ enthält, bleiben $\frac{\mathrm dV}{\mathrm dp}$ und $P$ nicht mehr gleichbedeutend; letztere Grösse stellt hier offenbar in Beziehung auf $p$ dasselbe vor, was $X^0$ in Beziehung auf $x$ im 15. Artikel; $\frac{\mathrm dV}{\mathrm dp}$ hingegen hat zwei verschiedene Werthe, nemlich $P-2\pi k$ und $P+2\pi k$, jenachdem $\mathrm dp$ als positiv oder als negativ betrachtet wird. Da nun $\int k\mathrm ds$ offenbar der ganzen auf die Fläche vertheilten Masse $M'$ gleich, und gemäss dem Lehrsatze des vorhergehenden Artikels $\int P\mathrm ds=2\pi M'$ wird, so hat man
\[
\int\frac{\mathrm dV}{\mathrm dp}.\mathrm ds=0\quad\text{oder}\quad\int\frac{\mathrm dV}{\mathrm dp}.\mathrm ds=4\pi M'
\]
jenachdem für $\frac{\mathrm dV}{\mathrm dp}$ der auf der innern, oder der auf der äussern Seite der Fläche geltende Werth überall verstanden wird, und es verhält sich also mit dem Integrale $\int\frac{\mathrm dV}{\mathrm dp}.\mathrm ds$ im erstern Falle genau eben so, als wenn die Masse $M'$ zum äussern Raume, im zweiten, als ob sie zum innern Raume gehörte.

Es gilt daher, bei irgendwie vertheilten Massen, die Gleichung
\[
\int\frac{\mathrm dV}{\mathrm dp}.\mathrm ds=4\pi M
\]
